\[ %%%%%%%%%%%%%%%%%%%%%%%%%%%%%%%%%%%%%%%%%%%%%%%%%%%%%%%%%%%%%%%%%%%%%%%%%%%%%%%% % PHYSGEN % A Physics-Guided Graph Neural Framework for Long-Horizon Dynamical System Prediction %%%%%%%%%%%%%%%%%%%%%%%%%%%%%%%%%%%%%%%%%%%%%%%%%%%%%%%%%%%%%%%%%%%%%%%%%%%%%%%% \title{ PHYSGEN: A Physics-Guided Graph Neural Framework for Long-Horizon Dynamical System Prediction } \author{ Your Name\\ Affiliation\\ Email } \date{} \] 

sections/00_abstract.tex
\[ \begin{abstract} Accurate long-horizon prediction of nonlinear dynamical systems remains a fundamental challenge in scientific machine learning. Although recent advances in graph neural networks, neural operators, and physics-informed learning have significantly improved predictive capabilities for complex physical processes, existing methods often suffer from cumulative error propagation, loss of physical consistency, and instability during autoregressive forecasting over extended temporal horizons. These limitations become particularly pronounced in highly nonlinear, multi-scale systems governed by coupled partial differential equations, where small prediction errors rapidly amplify and lead to physically implausible solutions. This paper introduces \textbf{PHYSGEN (Physics-Guided Graph Neural Framework)}, a unified computational framework designed for stable and physically consistent long-horizon prediction of dynamical systems. Unlike conventional graph neural architectures that primarily optimize data fidelity, PHYSGEN explicitly incorporates physical knowledge into every stage of the learning process through differentiable physics constraints, graph-based state propagation, conservation-aware regularization, and adaptive stability control. Within the proposed framework, complex physical systems are represented as dynamic attributed graphs whose nodes correspond to spatial discretization elements, particles, or finite-volume cells, while graph edges encode local physical interactions and information exchange. The temporal evolution of the system is modeled as iterative message passing constrained by governing physical principles, enabling simultaneous learning of latent dynamics and preservation of physically meaningful behavior throughout long prediction sequences. A central component of PHYSGEN is the \emph{Physics Guidance Module}, which continuously evaluates the consistency of predicted states with conservation laws, governing equations, and domain-specific physical constraints. Instead of enforcing hard analytical solutions, the framework integrates differentiable residual losses that guide optimization while maintaining flexibility for data-driven adaptation. This hybrid strategy effectively combines the expressive representation power of deep graph neural networks with the robustness of first-principles physical modeling. To improve long-term stability, PHYSGEN introduces an adaptive error accumulation controller that monitors prediction uncertainty and dynamically adjusts latent state propagation during autoregressive inference. This mechanism substantially reduces trajectory divergence, suppresses numerical instability, and improves robustness across chaotic and stiff dynamical regimes. The proposed framework is designed to be largely domain-independent and applicable to a broad spectrum of scientific and engineering problems, including computational fluid dynamics, thermal transport, reaction--diffusion systems, elasticity, climate simulation, plasma physics, and digital twins of complex physical infrastructures. Owing to its modular architecture, PHYSGEN naturally supports heterogeneous graph representations, multi-scale simulations, and differentiable scientific computing pipelines. Theoretical analysis demonstrates that integrating graph-based message passing with continuous physics-guided regularization significantly improves long-horizon stability while preserving expressive learning capacity. The proposed architecture establishes a general computational paradigm for physics-guided graph intelligence that bridges modern geometric deep learning and scientific machine learning. \end{abstract} \textbf{Keywords:} Scientific Machine Learning, Physics-Guided Learning, Graph Neural Networks, Long-Horizon Prediction, Scientific Computing, Neural Operators, Digital Twins, Computational Physics, Differentiable Simulation, Physics-Informed Artificial Intelligence. \] 
